\subsection{Créditos en estudios avanzados}
Bandeja y expediente del trámite.
\captura{creditos_titulacion_02_bandeja}{Bandeja de validación de trámites.}
\captura{creditos_titulacion_01_expediente_ronda0}{Expediente del trámite con documentos por validar.}
\subsubsection*{Solicitud de la opción}
Valida o devuelve con observaciones cada documento; coteja el original cuando se indique.
\begin{longtable}{p{0.55\linewidth}p{0.35\linewidth}}
\toprule
\textbf{Documento} & \textbf{Observaciones}\\
\midrule
\endhead
Solicitud firmada por el pasante & La plataforma lo genera. Debe estar fechada antes de iniciar los estudios avanzados.\\
Certificado de estudios de licenciatura & \\
Certificado de servicio social & \\
Plan de estudios oficial de los estudios avanzados & Debe avalar cursos, créditos, actividades con valor curricular, plazos mínimos, perfiles y contenidos programáticos.\\
Registro del programa ante la Dirección General de Profesiones (o validez oficial) & Solo para estudios que no se realicen en la UAEMéx.\\
\bottomrule
\end{longtable}

Dictamen de afinidad y autorización. La Coordinación de Docencia dictamina la afinidad y la Subdirección Académica autoriza.
Datos que registras:
\begin{itemize}
  \item Fecha de autorización
\end{itemize}
Se genera: Carta de afinidad, que firma la Subdirección.
\subsubsection*{Estudios avanzados cursados}
Valida o devuelve con observaciones cada documento; coteja el original cuando se indique.
\begin{longtable}{p{0.55\linewidth}p{0.35\linewidth}}
\toprule
\textbf{Documento} & \textbf{Observaciones}\\
\midrule
\endhead
Carta de aceptación a los estudios avanzados & Con fecha posterior a la carta de afinidad.\\
Dictamen de equivalencia de la Secretaría de Investigación y Estudios Avanzados & \\
Certificado de estudios y constancia de calificaciones de créditos cubiertos & Especifica que se cubrieron en los plazos mínimos.\\
\bottomrule
\end{longtable}
\subsubsection*{Documentación para la toma de protesta}
Valida o devuelve con observaciones cada documento; coteja el original cuando se indique.
\begin{longtable}{p{0.55\linewidth}p{0.35\linewidth}}
\toprule
\textbf{Documento} & \textbf{Observaciones}\\
\midrule
\endhead
Certificado de estudios de bachillerato & \\
Constancia de autenticación de bachillerato & Se tramita en la escuela donde cursaste el bachillerato. Solo si tus estudios de bachillerato no fueron en instituciones dependientes o incorporadas a la UAEM.\\
Acta de nacimiento & Copia certificada vigente (del año en curso).\\
CURP & Debe mostrar la leyenda «Validada por el Registro Civil».\\
Constancia o certificado de aprobación del idioma inglés & Solo planes 98 en adelante, si en tu certificado no aparecen Inglés C1 e Inglés C2.\\
Constancia de no adeudo a la Biblioteca de la Facultad & El recibo de pago se tramita en Control Escolar de la Facultad; la constancia, en la Biblioteca.\\
Constancia de no adeudo a la Biblioteca Central & \\
Contra-recibo de la donación de libros a la Biblioteca de la Facultad & \\
Constancia de no adeudo a la Dirección de Control Escolar de la UAEM & Se solicita en el Portal de alumnos: controlescolar.uaemex.mx.\\
Recibo de no adeudo a la UAEM & Se solicita en el Portal de alumnos de Control Escolar.\\
Formato de datos personales (para el título) & La plataforma lo genera. La plataforma lo llena con tus datos; lo firmas electrónicamente.\\
Comprobante del sistema de Seguimiento de Egresados & Se obtiene en alumniversitario.uaemex.mx/1dae.php.\\
Dictamen de autorización de los HH. Consejos & La plataforma lo genera. Se solicita en la Subdirección Académica con el formato del Art. 5. Solo si excediste el plazo de 10 años para presentar la evaluación profesional (Art. 5, REP).\\
Recibo de pago: derechos de evaluación profesional (\$230) & Se solicita en Control Escolar de la Facultad.\\
Recibo de pago: expedición de título (\$1,187) & Se solicita en Control Escolar de la Facultad.\\
\bottomrule
\end{longtable}
\subsubsection*{Fecha, hora y lugar de la toma de protesta}

Fecha y lugar de la toma de protesta. Titulación fija la fecha. En esta opción la comisión es siempre el Director (presidente) y el Subdirector Académico (secretario); el acta oficial se genera en otro sistema y se firma en papel.
Datos que registras:
\begin{itemize}
  \item Fecha
  \item Hora
  \item Lugar
\end{itemize}
\subsubsection*{Resultado de la evaluación}

Registrar resultado. 
Datos que registras:
\begin{itemize}
  \item Resultado
  \item Fecha de la evaluación
  \item Número de acta (de evaluación profesional; si se aplaza, de aplazamiento)
\end{itemize}
\subsubsection*{Verificación final del expediente}
Integración De Expediente: Créditos En Estudios Avanzados.

Documentos que verifica el responsable, en el orden del formato oficial:
\begin{enumerate}
  \item Dos fotografías tamaño semitítulo, blanco y negro, papel mate, fondo gris claro, de frente, con retoque, varones con camisa blanca, saco y corbata oscuros, mujeres con blazer oscuro --- \emph{físico}
  \item Acta de evaluación profesional --- \emph{físico}
  \item Copia del certificado de bachillerato con firma y sello de cotejo con el original --- \emph{plataforma}
  \item Constancia de autenticación si el bachillerato no es UAEM --- \emph{plataforma (sólo si aplica)}
  \item Dictamen de revalidación expedido por la Secretaría de Educación Pública (sólo si aplica) --- \emph{físico (sólo si aplica)}
  \item Acta de nacimiento --- \emph{plataforma}
  \item Certificado total de licenciatura --- \emph{plataforma}
  \item Solicitud por escrito del egresado dirigida a la Subdirección Académica con firma y sello de la misma --- \emph{plataforma}
  \item Certificado que ampare el 75 \% de maestría o doctorado integral o el 100 \% del diplomado superior --- \emph{plataforma}
  \item Dictamen de equivalencia expedido por la Secretaría de Ciencia para cualquier diplomado superior universitario, maestría o doctorado integral --- \emph{plataforma}
  \item Autorización para continuar con trámites de titulación con fundamento en el artículo 5° del Reglamento de Evaluación Profesional (sólo si aplica) --- \emph{plataforma (sólo si aplica)}
  \item Constancia de no adeudo a Servicios Escolares --- \emph{plataforma}
  \item Formato de revisión para autorización de evaluación profesional --- \emph{plataforma}
  \item Formato de integración de expediente --- \emph{plataforma}
  \item Formato de datos del interesado para expedición de título profesional --- \emph{plataforma}
\end{enumerate}
\captura{creditos_titulacion_verificacion_01}{Lista de verificación final con estado de cada renglón.}
\captura{creditos_titulacion_verificacion_02}{Expediente verificado y exportación para Servicios Escolares.}
Marca con «Entregado y revisado» lo que se entrega en físico y con «No aplica» lo condicional; con todo cubierto, pulsa «Hago constar que el expediente está completo y verificado». Después, en \emph{Expedientes a Servicios Escolares}, elige los verificados y pulsa «Exportar carpeta comprimida para Servicios Escolares».
